% Generated by supplementary/check_manuscript.py from verified model data.
\begin{table}[htbp]
\centering\small
\begin{tabular}{@{}clcl@{}}
\toprule
$\mu$ & A unitary realization & $D$ & Reference \\
\midrule
$1$ & $\mathrm{Vec}_{C_2\times C_2}$ & $4$ & \cite[\S3.4]{LPR} \\
$1$ & $\mathrm{Fib}\boxtimes\mathrm{Vec}_{C_2}$ & $5+\sqrt{5}$ & \cite[\S3.4]{LPR} \\
$1$ & $\operatorname{Rep}(D_{10})$ & $10$ & \cite[\S3.4]{LPR} \\
$1$ & $\operatorname{Ad}(\mathrm{SU}(2)_6)$ & $8+4\sqrt{2}$ & \cite[Theorem 3.1]{EIP} \\
$1$ & $\mathrm{Fib}\boxtimes\mathrm{Fib}$ & $(15+5\sqrt{5})/2$ & \cite[\S3.4]{LPR} \\
$1$ & $\operatorname{Ad}(\mathrm{SU}(2)_7)$ & $\delta_7$ & \cite[\S3.4]{LPR} \\
$1$ & $\mathrm{Vec}_{C_4}$ & $4$ & \cite[\S3.4]{LPR} \\
$1$ & Tambara--Yamagami on $C_3$ & $6$ & \cite{TambaraYamagami} \\
$1$ & Principal even part of $S'$ & $8+4\sqrt{2}$ & \cite[Theorem 1]{LMP} \\
$2$ & Principal even part of $2D2$ & $20+8\sqrt{5}$ & \cite[Theorem 3.1]{EIP} \\
$2$ & Haagerup category $\mathcal H_1$ & $(39+9\sqrt{13})/2$ & \cite[Table 2]{GrossmanSnyder} \\
$2$ & $\operatorname{Rep}(A_4)$ & $12$ & \cite[Theorem 1.1]{Larson} \\
$3$ & $\operatorname{Ad}(E_8)$ & $\delta_E$ & \cite{IzumiE8,EvansJones} \\
$3$ & Near-group type $(C_3,3)$ & $(21+3\sqrt{21})/2$ & \cite[Example 9.4]{Izumi} \\
$6$ & Near-group type $(C_3,6)$ & $24+12\sqrt{3}$ & \cite[Theorem 10.19]{Izumi} \\
\bottomrule
\end{tabular}
\caption{One unitary realization of each surviving based ring, with its
multiplicity $\mu$ and global Frobenius--Perron dimension $D$.
Here $D_{10}$ has order ten, $S'$ has index $3+2\sqrt2$, and the
near-group category of type $(C_3,3)$ is the Izumi--Xu category.}
\label{finite:tab:models}
\smallskip
\begin{minipage}{\textwidth}\footnotesize
The numbers $\delta_7$ and $\delta_E$ are the largest real roots of
$x^3-27x^2+162x-243$ and
$x^4-45x^3+375x^2-1125x+1125$, respectively.
The supplementary model check verifies these polynomials, the positive
dimension characters, and all tensor identifications exactly.
\end{minipage}
\end{table}
